\chapter{Structures, Satisfaction, and Translation}
\label{ch:fol-semantics}

\section{Structures and variable assignments}
\label{sec:fol-semantics-01}

\begin{definition}[First-order structure]
Let \(\Lang\) be a signature.  A \emph{\(\Lang\)-structure}
\(\Meta\) consists of:
\begin{enumerate}
  \item a nonempty set \(\dom\Meta\), called its domain;
  \item for each \(n\)-ary function symbol \(f\), a function
        \(f^{\Meta}:(\dom\Meta)^{n}\to\dom\Meta\); and
  \item for each \(n\)-ary relation symbol \(R\), a relation
        \(R^{\Meta}\subseteq(\dom\Meta)^{n}\).
\end{enumerate}
For a constant symbol \(c\), the zero-ary function is identified with its
value \(c^{\Meta}\in\dom\Meta\).  The equality sign is always interpreted
as actual identity on \(\dom\Meta\).
\end{definition}

\begin{definition}[Signature extension, expansion, and reduct]
Let \(\Lang=\tuple{\mathcal F,\mathcal R,\arity_{\mathcal F},
\arity_{\mathcal R}}\) and
\(\Lang'=\tuple{\mathcal F',\mathcal R',\arity_{\mathcal F'},
\arity_{\mathcal R'}}\).  The signature \(\Lang'\) \emph{extends}
\(\Lang\) iff \(\mathcal F\subseteq\mathcal F'\),
\(\mathcal R\subseteq\mathcal R'\), and both arity maps of \(\Lang'\)
restrict to the corresponding arity maps of \(\Lang\).

If \(\Meta\) is a \(\Lang\)-structure, a \(\Lang'\)-structure
\(\Meta'\) is an \emph{expansion} of \(\Meta\) iff it has the same domain
and gives every symbol of \(\Lang\) the same interpretation.  If
\(\Meta'\) is a \(\Lang'\)-structure, its \emph{reduct} to \(\Lang\),
written \(\Meta'\restr\Lang\), is the \(\Lang\)-structure with the same
domain and the interpretations of the symbols in \(\Lang\).  Expansion adds
interpretations and reduct forgets them; neither operation changes the domain,
so neither term means substructure.  These operations are required in the
Henkin, L\"owenheim--Skolem, and Skolemization arguments below.
\end{definition}

\begin{definition}[Variable assignment and update]
A \emph{variable assignment} in \(\Meta\) is a function
\(s:\Var\to\dom\Meta\).  For \(a\in\dom\Meta\), define the updated
assignment \(s[x\mapsto a]\) by
\[
 s[x\mapsto a](y)=
 \begin{cases}
   a,&y=x,\\
   s(y),&y\ne x.
 \end{cases}
\]
\end{definition}

\begin{definition}[Term value]
The \emph{value} of a term \(t\) in \(\Meta\) under \(s\), written
\(\sem{t}^{\Meta}_s\), is defined by
\[
 \sem{x}^{\Meta}_s=s(x),
 \qquad
 \sem{f(t_1,\ldots,t_n)}^{\Meta}_s
 =f^{\Meta}(\sem{t_1}^{\Meta}_s,\ldots,\sem{t_n}^{\Meta}_s).
\]
\end{definition}

\begin{lemma}[Term coincidence]
If assignments \(s\) and \(s'\) agree on \(\operatorname{Var}(t)\), then
\(\sem{t}^{\Meta}_s=\sem{t}^{\Meta}_{s'}\).
\end{lemma}
\begin{proof}
Induct on the construction of \(t\).  The variable case is the hypothesis.
In a function case, apply the induction hypothesis to each argument and then
apply the function \(f^{\Meta}\).
\end{proof}

\begin{lemma}[Term-substitution lemma]
\label{lem:term-substitution}
For all terms \(r,t\), assignments \(s\), and variables \(x\),
\[
 \sem{r[t/x]}^{\Meta}_s
 =\sem{r}^{\Meta}_{s[x\mapsto\sem{t}^{\Meta}_s]}.
\]
\end{lemma}
\begin{proof}
Induct on \(r\).  If \(r=x\), both sides equal
\(\sem{t}^{\Meta}_s\).  If \(r=y\ne x\), both sides equal \(s(y)\).  The
function case follows componentwise from the induction hypotheses and the
definition of term value.
\end{proof}

